\begin{ouverture}[college2]{9}{Interactions, forces et poids}{Partie 2 · Mouvements et interactions}{Comment représenter les actions sur un objet ?}
\objectif{Identifier les interactions de contact et à distance.}
\objectif{Représenter une force à une échelle donnée.}
\objectif{Distinguer masse et poids et exploiter P = mg.}
\prerequis{Masse en kg et vitesse.}
\prerequis{Directions et proportionnalité.}
\end{ouverture}
\pcPage{Le cours}{Interactions, forces et poids}
\phantomsection\label{ess-9}
\essSection{1}{Choisir le système}
Le système est l’objet ou l’ensemble d’objets que l’on étudie. Un diagramme objet-interactions relie ce système à son environnement : support, fil, air, Terre… Une interaction de contact exige un contact matériel ; la gravitation agit à distance. Le diagramme recense les partenaires ; il ne donne pas à lui seul la valeur des forces.\par
\pcFigure{figures/b09-doi.pdf}{Une balle suspendue : interactions avec le fil et la Terre ; l’air est négligé dans ce modèle.}
\essSection{2}{Modéliser par une force}
Une force représente l’action d’un objet sur un autre. On la caractérise par un point d’application, une direction, un sens et une valeur en newtons, N. La flèche part du point choisi ; sa longueur respecte une échelle annoncée. Direction verticale et sens vers le bas sont deux informations différentes. Un dynamomètre mesure une force selon son axe, après contrôle du zéro.\par
\clearpage
\essSection{3}{Poids et masse}
La masse m caractérise la quantité de matière dans les situations étudiées et se mesure en kg avec une balance. Le poids P est une force de pesanteur, en N, mesurée avec un dynamomètre. $P=mg$, avec m en kg et g en N/kg. Près de la surface terrestre, on prend souvent $g=9{,}8$ ou $10\unite{N/kg}$ selon l’énoncé. Le poids est dirigé vers le centre de la Terre dans le modèle local.\par
\pcFigure{figures/b09-poids.pdf}{Même masse, poids différents : les valeurs de g dépendent du lieu.}
\essSection{4}{Effets et équilibre}
Une action peut déformer un objet ou modifier son mouvement, en valeur ou en direction. Dans un modèle d’objet immobile soumis à deux forces qui se compensent, elles ont même direction, sens opposés et même valeur. Elles s’exercent sur le même système. Deux forces d’une interaction entre deux objets s’exercent, elles, sur des systèmes différents : on ne les additionne pas sur un seul objet.\par
\begin{pcEncadre}{method}{Un exemple expliqué}Une trousse de 350 g est immobile sur une table. $m=0{,}350\unite{kg}$, donc $P=0{,}350\times10=3{,}5\unite N$. Dans le modèle à deux forces, le support exerce 3,5 N vers le haut. À 1 cm pour 1 N, chaque flèche mesure 3,5 cm.\end{pcEncadre}
\pcPage{Expérimenter et retenir}{Interactions, forces et poids}
\phantomsection\label{activite-9}
\begin{pcEncadre}{method}{Investigation en classe}Mesurer le poids de plusieurs masses avec un dynamomètre adapté. Tracer P en fonction de m en kg, puis calculer P/m hors zéro. Comparer la pente à la valeur locale de g. Dessiner ensuite les forces sur une masse suspendue immobile, en distinguant les actions sur la masse et sur l’appareil.\end{pcEncadre}
\pcSousTitre{Ce que je dois savoir expliquer}\begin{questions}
\item Nommer l’auteur et le receveur d’une force.
\item Décrire les quatre caractéristiques d’une flèche.
\item Convertir les grammes en kilogrammes avant P = mg.
\end{questions}
\begin{pcEncadre}{attention}{Une erreur à éviter}Un objet posé sur une table a un poids. Son immobilité ne signifie pas absence de forces : dans le modèle, l’action du support compense le poids.\end{pcEncadre}
\pcRepere{cor-activite-9}{Pistes pour l’investigation}\pcGrille
\clearpage\markright{Je démarre}
\rubrique{Je démarre}
\par\noindent\begin{minipage}{\linewidth}
\begin{exercice}[1]{Contact ou distance}
\begin{questions}
\item Classer l’action d’un fil sur une masse suspendue.
\item Classer celle de la Terre sur cette masse.
\end{questions}
\end{exercice}
\end{minipage}
\par\noindent\begin{minipage}{\linewidth}
\begin{exercice}[1]{Deux appareils}
\begin{questions}
\item Quel appareil mesure une masse ?
\item Quel appareil mesure une force ?
\end{questions}
\end{exercice}
\end{minipage}
\par\noindent\begin{minipage}{\linewidth}
\begin{exercice}[1]{Deux unités}
\begin{questions}
\item Donner l’unité SI de la masse.
\item Donner celle du poids.
\end{questions}
\end{exercice}
\end{minipage}
\par\noindent\begin{minipage}{\linewidth}
\begin{exercice}[1]{Direction et sens}
\begin{questions}
\item Le poids est-il « vers le bas » ou « vertical » : quelle information donne chaque expression ?
\item Citer les deux autres caractéristiques de sa représentation.
\end{questions}
\end{exercice}
\end{minipage}
\par\noindent\begin{minipage}{\linewidth}
\begin{exercice}[1]{Un premier calcul}
\begin{questions}
\item Calculer le poids de 2 kg pour $g=10\unite{N/kg}$.
\item Une masse de 200 g vaut-elle 200 kg ?
\end{questions}
\end{exercice}
\end{minipage}
\clearpage\markright{Je m’entraîne}
\rubrique{Je m’entraîne}
\par\noindent\begin{minipage}{\linewidth}
\begin{exercice}[2]{Une petite masse}
\begin{questions}
\item Calculer P pour 150 g et $g=9{,}8\unite{N/kg}$.
\item Donner une valeur arrondie au dixième de N.
\end{questions}
\end{exercice}
\end{minipage}
\par\noindent\begin{minipage}{\linewidth}
\begin{exercice}[2]{Retrouver la masse}
\begin{questions}
\item Un objet pèse 24 N pour $g=10\unite{N/kg}$. Déterminer sa masse.
\item L’exprimer en g.
\end{questions}
\end{exercice}
\end{minipage}
\par\noindent\begin{minipage}{\linewidth}
\begin{exercice}[2]{Changer de lieu}
\begin{questions}
\item Une masse de 60 kg est transportée sur la Lune, avec $g=1{,}6\unite{N/kg}$. Donner masse et poids.
\item Comparer à son poids terrestre avec $g=10\unite{N/kg}$.
\end{questions}
\end{exercice}
\end{minipage}
\par\noindent\begin{minipage}{\linewidth}
\begin{exercice}[2]{Une échelle}
\begin{questions}
\item Une force de 12 N est représentée à 1 cm pour 4 N. Calculer la longueur.
\item Une flèche de 2,5 cm représente quelle valeur à cette échelle ?
\end{questions}
\end{exercice}
\end{minipage}
\par\noindent\begin{minipage}{\linewidth}
\begin{exercice}[2]{Le livre immobile}
\begin{questions}
\item Nommer les deux partenaires principaux d’un livre immobile sur une table horizontale.
\item Décrire leurs forces dans le modèle à deux forces.
\end{questions}
\end{exercice}
\end{minipage}
\clearpage\markright{Je raisonne}
\rubrique{Je raisonne}
\par\noindent\begin{minipage}{\linewidth}
\begin{exercice}[3]{Une pente mesurée}
\pcFigureEx{figures/rev-ex-poids-graphe.pdf}{}
\begin{questions}
\item On mesure 1,0 N pour 0,10 kg et 2,0 N pour 0,20 kg. Estimer g.
\item Que suggère un graphe linéaire passant près de l’origine ?
\end{questions}
\end{exercice}
\end{minipage}
\par\noindent\begin{minipage}{\linewidth}
\begin{exercice}[3]{Une mesure négative}
\begin{questions}
\item Un dynamomètre indique −0,2 N sans charge. Quelle vérification faire ?
\item Pourquoi ce décalage gêne-t-il surtout une petite mesure ?
\end{questions}
\end{exercice}
\end{minipage}
\par\noindent\begin{minipage}{\linewidth}
\begin{exercice}[3]{Des forces confondues}
\begin{questions}
\item Un élève dessine sur le livre son poids et la force qu’il exerce sur la table. Quelle erreur commet-il ?
\item Quelle force doit remplacer la seconde ?
\end{questions}
\end{exercice}
\end{minipage}
\par\noindent\begin{minipage}{\linewidth}
\begin{exercice}[3]{Un virage}
\begin{questions}
\item Une voiture garde une même valeur de vitesse dans un virage. Son mouvement est-il inchangé ?
\item Une force peut-elle modifier seulement la direction de la vitesse ?
\end{questions}
\end{exercice}
\end{minipage}
\par\noindent\begin{minipage}{\linewidth}
\begin{exercice}[3]{Une affirmation à discuter}
\begin{questions}
\item « Une force de 5 N suffit toujours à déplacer n’importe quel objet. » Pourquoi faut-il connaître les autres actions ?
\item Donner un cas où cette force est compensée.
\end{questions}
\end{exercice}
\end{minipage}
\pcPage{Défi du chapitre}{Interactions, forces et poids}
\phantomsection\label{defi-ex-9}
\pcDefiTitre{9}{La caisse suspendue}
\begin{questions}
\item Une caisse de 4,0 kg est suspendue immobile, $g=9{,}8\unite{N/kg}$. Calculer P.
\item Tracer les forces qui s’appliquent à la caisse, à l’échelle 1 cm pour 10 N.
\item Pourquoi ces forces de compensation ne sont-elles pas les deux forces d’une même interaction ?
\end{questions}
\pcGrille
